% !PRINTED 69
\weberpage{69}
This function $\varrho$ then belongs to a certain group $\mathfrak{A}$,
which is a divisor of $\mathfrak{N}$. We also call the group $\mathfrak{A}$ and
the field $\Omega(\varrho)$ \gesperrt{belonging together}, and from this we draw the
theorem:

\gesperrt{1. To every divisor $\mathfrak{A}$ of $\mathfrak{N}$ there belongs a certain
cyclotomic field $\Omega(\varrho)$ contained in $\Omega_m$, and
conversely to every divisor $\Omega(\varrho)$ of $\Omega_m$
there belongs a certain divisor $\mathfrak{A}$ of $\mathfrak{N}$ as group, in the
sense that, if $a$ is a number from $\mathfrak{A}$, all
numbers of the field $\Omega(\varrho)$ admit the permutations $(r,r^a)$,
and that conversely every number in $\Omega_m$
that admits these permutations is contained in $\Omega(\varrho)$.}

The Galois group of such a field $\Omega(\varrho)$ we obtain, according to Bd. I, §. 156, by decomposing $\mathfrak{N}$ into the system of
subsidiary groups:
\[
\mathfrak{N} = \mathfrak{A} + \mathfrak{A}_1 + \mathfrak{A}_2 + \cdots,
\]
and seeking the permutations produced among the subsidiary groups $\mathfrak{A}, \mathfrak{A}_1, \mathfrak{A}_2 \ldots$ by composition
with the elements of $\mathfrak{N}$. But this group is, according to §. 4, isomorphic with the
group $\mathfrak{N}/\mathfrak{A}$, and hence also isomorphic with the group reciprocal to $\mathfrak{A}$
(§. 12), which we shall denote by $\mathfrak{B}$. If $a$ is the degree of $\mathfrak{A}$
and $b$ that of $\mathfrak{B}$, then $ab = \nu$, and $\varrho$ satisfies an irreducible
Abelian equation of degree $b$.

\gesperrt{2. We thus obtain all cyclotomic fields
if for every modulus $m$ we form all the
divisors $\mathfrak{A}$ of the group $\mathfrak{N}$, seek for each of these
groups $\mathfrak{A}$ an associated function $\varrho$,
and from these derive the fields $\Omega(\varrho)$.}

But there remains the question whether in this process one and
the same field $\Omega$ can occur more than once, whereby we are led to
the investigation of the common divisors of two cyclotomic fields.

According to the definition given above, one must carefully distinguish between the group to which a field $\Omega$ \gesperrt{belongs}
and the Galois group of the field; the two groups are
reciprocal to each other. Thus the field $\Omega_m$ itself belongs to the
unit group, while its Galois group is $\mathfrak{N}$,
